\documentclass[12pt,table, hidelinks]{extarticle} % что то меняется , не понятно article or extarticle...
\usepackage[T1, T2A]{fontenc}
\usepackage[utf8]{inputenc}
\usepackage[english,russian]{babel}
%\usepackage[dvips]{graphicx}
\usepackage{graphicx}
%\usepackage{svg}


\usepackage{unicode-math} % Задать шрифт для математических формул
\setmathfont{Stix Two Math}
%\setmathfont{Arial}

% для KOMA script чтобы в документе менять формат А3 А4 и ориентацию
\usepackage{typearea}

\usepackage[
a4paper,
includehead,
includefoot,
nomarginpar,% We don't want any margin paragraphs
%textwidth=10cm,% Set \textwidth to 10cm
headheight=10mm,% Set \headheight to 10mm
top=1mm,
bottom=7mm, 
left=22mm,
right=7mm,
]{geometry}

%\savegeometry{standard}
\setlength{\topmargin}{-27mm}
\setlength{\headsep}{7mm} % Устанавливаем отступ между колонтитулом и текстом 10pt !!! Сдвигает а4 landscape
%\setlength{\footskip}{10pt} % Устанавливаем отступ между текстом и колонтитулом внизу страницы 20pt
%\setlength{\topskip}{0pt} % Установка вертикального промежутка между верхним краем текста и колонтитулом на 0pt
\setlength{\textheight}{268mm}
\setlength{\footskip}{-5mm}


\usepackage{fancyhdr}

%\usepackage{lmodern}


%\usepackage{longtable}
%\usepackage{amsfonts}

\usepackage{xcolor} % ПРИ КОМПИЛЯЦИИ В ЛИНУКС УБРАТЬ ЭТУ СТРОКУ, да вообще ее убрать
\usepackage{colortbl}

%\usepackage{array}

\usepackage{hyperref}

\graphicspath{{./images/}}

\usepackage{fontspec}
\setsansfont[
        Ligatures=TeX,
        BoldFont={Arial Narrow Bold},
        ItalicFont={Arial Narrow Italic},
        BoldItalicFont={Arial Narrow Bold Italic}
    ]{Arial Narrow}
\renewcommand\familydefault{\sfdefault}
%\setmainfont{Arial Narrow}
\usepackage{fancyhdr}

\usepackage{parskip}% для абзацев и отступа
\setlength{\parindent}{1cm}
\setlength{\parskip}{1mm}
\setlength{\intextsep}{0mm} % отступ между плавающими элементами


\usepackage{indentfirst}

% Определяем цвета фирмы
\definecolor{unired}{RGB}{247, 148, 29} 
\definecolor{uniblue}{RGB}{30, 73, 115}
\definecolor{unigray}{RGB}{88, 89, 91}

\usepackage{longtable,booktabs}
\usepackage{calc} % for calculating minipage widths


% Это для таблицы, чтобы жестко ширину столбцов выставлять НЕ ТРЕБУЕТСЯ???
\usepackage{array}
\newcolumntype{L}[1]{>{\raggedright\let\newline\\\arraybackslash\hspace{0pt}}m{#1}}
\newcolumntype{C}[1]{>{\centering\let\newline\\\arraybackslash\hspace{0pt}}m{#1}}
\newcolumntype{R}[1]{>{\raggedleft\let\newline\\\arraybackslash\hspace{0pt}}m{#1}}


\usepackage{bigstrut}

\usepackage[labelsep=endash]{caption} % выравнить надписи к рисункам тире вместо двоеточия
\DeclareCaptionLabelFormat{myformat}{\textls[75]{#1 #2}} % Здесь устанавливается разрежение текста для надписей Рисунок и Таблица
\captionsetup[figure]{name=Рисунок, justification=centering, labelformat=myformat} 
\captionsetup[table]{justification=justified, labelformat=myformat} % Убираем отступ при переносе строки в наименовании таблицы
\setlength{\belowcaptionskip}{0mm} % отступ после подписи
%\setlength{\abovecaptionskip}{0mm} % отступ до подписи

% нумерация таблиц и рисунков ИСПОЛЬЗУЕТСЯ!!!
\counterwithin*{figure}{subsection}
\counterwithin*{table}{subsection}
% * means don't change the representation, because we'll do it now

\usepackage{titlesec}

% Форматирование заголовков секций и подразделов
\titleformat{\section}
  {\normalfont\Large\bfseries\MakeUppercase}{\thesection}{1em}{}
\titleformat{\subsection}
  {\normalfont\large\bfseries\MakeUppercase}{\thesubsection}{1em}{}

\renewcommand{\thefigure}{%
  \ifnum\value{subsection}>0
    \thesubsection.%
  \else
    \thesection.%
  \fi
  \arabic{figure}%
}
\renewcommand{\thetable}{%
  \ifnum\value{subsection}>0
    \thesubsection.%
  \else
    \thesection.%
  \fi
  \arabic{table}%
}
%%%%%%%%%%%%%%% ЭТОТ КОД ДЕЛАЕТ ТОЖЕ САМОЕ? %%%%%%%%%%%%%%%%%%%
%\usepackage{chngcntr}
%\counterwithin{figure}{subsection}
%\counterwithin{table}{subsection}
%\renewcommand{\thefigure}{\thesubsection.\arabic{figure}}
%\renewcommand{\thetable}{\thesubsection.\arabic{table}}
%%%%%%%%%%%%%%%%%%%%%%%%%%%%%%%%%%%%%%%%%%%%%%%%%%%%%%%%%%%%%%%

\usepackage{tabularx}
\usepackage[table]{xcolor}
\usepackage{multirow} % Для центрирования содержимого ячейки по вертикали - ИСПОЛЬЗУЕТСЯ!!!
\usepackage{amsmath} % для вывода формул
\usepackage{eqexpl} % для вывода пояснений к формулам
\eqexplSetDelim{--}
\eqexplSetIntro{где}

\usepackage{titlesec}
\usepackage{appendix}
%\titlespacing{\chapter}{0pt}{-30pt}{12pt} как пример
\titlespacing{\section}{0pt}{0pt}{5mm}
\titlespacing{\subsection}{0pt}{3mm}{0mm}
\titlespacing{\subsubsection}{0pt}{2mm}{2mm}


%\usepackage[bitstream-charter]{mathdesign} % проверить когда формулы будут в тексте
\uchyph=0 % отключаем перенос в словах с заглавной буквы - чтобы ГЗоткл не переносил например


% Заменяем точки на тире в перечислении (списке)
% - дефис, -- тире, --- длинное тире
\usepackage{enumitem}

\setlist[itemize]{label=---}
%\setlist[itemize]{label=$\circ$} % окружность
%%%%%%%%%%%%%%%%%%%%%%%%
\usepackage{float}


% не допускает в конце странцы заголовок и одна или две строки , переносит на след страницу
% необходимо указывать непосредственно в разделе
\usepackage{needspace} %ИСПОЛЬЗУЕТСЯ!!!
%%%%%%%%%%%%%%%%%%%%%%%%

\usepackage{tabularray}
\usepackage{tabularx}


\usepackage{tocloft}  % точки для глав в оглавлении ИСПОЛЬЗУЕТСЯ!!!
\usepackage{tocvsec2} % для подавления отображения подразделов в toc для приложений, не знаю, если подразделы появятся
\renewcommand{\cftsecleader}{\cftdotfill{\cftdotsep}} % точки для глав в оглавлении ИСПОЛЬЗУЕТСЯ!!!

\usepackage{pdflscape}
\titleformat{\section}{\normalfont\large\bfseries}{\thesection}{1em}{}

%\usepackage{lscape} % УДАЛИТЬ???
\usepackage{lipsum}
\usepackage[absolute]{textpos}
\usepackage{etoolbox}

\usepackage{pdfpages} % попробуем вставить векторный рисунок


%\titlespacing{\subsection}{0pt}{\parskip}{-\parskip} % отступ между сабсекцией и лонгтейбл?
%\titlespacing{command}{left spacing}{before spacing}{after spacing}[right]

%\titlespacing\section{0pt}{12pt plus 4pt minus 2pt}{0pt plus 2pt minus 2pt}
%\titlespacing\subsection{0pt}{12pt plus 4pt minus 2pt}{0pt plus 2pt minus 2pt}
%\titlespacing\subsubsection{0pt}{12pt plus 4pt minus 2pt}{0pt plus 2pt minus 2pt}

%\titlespacing{\section}{0pt}{\parskip}{-\parskip}
%\titlespacing{\subsection}{0pt}{\parskip}{-\parskip}
%\titlespacing{\subsubsection}{0pt}{\parskip}{-\parskip}

\usepackage{fontawesome} % Подключить пакет шрифтов fontawesome

\usepackage{microtype} % для разрежения текста

\usepackage{etoolbox} % для проверки логических условий


